\chapter{Канонический потенциал и нерадиальная устойчивость дефекта}

Радиальная проверка нашла гладкое решение и положительную вторую вариацию
внутри однопараметрического сектора. Перед расширением на все пять компонент
$Q$ необходимо, однако, вернуться к уже выведенному фазовому функционалу.
Иначе нерадиальный гессиан зависел бы от произвольно продолженного объёмного
потенциала.

\section{Литературная и проектная граница}

В теории Ландау--де Жена устойчивость радиального ежа относительно
произвольных $Q$-тензорных возмущений является отдельной задачей. В
\path{arXiv:1404.1729} доказана устойчивость плавящегося ежа вблизи
критической температуры и её потеря в глубокой нематической фазе. Тем самым
сама пятикомпонентность поля способна менять ответ, даже когда радиальный
профиль существует. Монотонность и единственность радиального минимизатора
не заменяют эту проверку; см. \path{arXiv:1212.1072}. Наконец, решёточный
анализ $B=2$-ежа Скирма в \path{arXiv:hep-ph/9509421} показывает общий
механизм: нерадиальная отрицательная мода может описывать распад
симметричной конфигурации на несколько солитонов.

В нашем проекте имеется дополнительное требование. Объёмный потенциал уже
получен как каноническая свободная энергия проекторной фазы, поэтому новая
проверка обязана использовать именно её, а не только симметрийно допустимый
квартичный заменитель.

\section{Исправление области действия предыдущего результата}

Пусть $R=I_3/3+Q$ --- положительная матрица следа один. Канонический
локальный функционал Тома VI равен
\begin{equation}
 \mathcal F_\beta(R)=
 \Tr(R\log R)+\beta\left[
 \frac27\left(1-\frac{\Tr(R^2)^2}{\Tr(R^3)}\right)
 +1-\Tr(R^2)\right].
 \label{eq:v6-canonical-local-free-energy}
\end{equation}
При $\beta=\beta_c=1.542669540860\ldots$ изотропное состояние $I_3/3$ и
упорядоченное состояние $R_*$ имеют одинаковую свободную энергию. Поэтому
родительский объёмный член должен быть
\begin{equation}
 V_{\rm can}(Q)=
 \mathcal F_{\beta_c}(I_3/3+Q)-\mathcal F_{\beta_c}(R_*).
 \label{eq:v6-canonical-defect-potential}
\end{equation}

В предыдущей главе использовалось допустимое квартичное продолжение
\begin{equation}
 V_4(u)=\Delta^4(1-u^2)^2.
\end{equation}
Оно положительно при $u=0$, тогда как
$V_{\rm can}(0)=V_{\rm can}(Q_*)=0$. Следовательно, прежние числа
$r_{1/2}=1.03944\ldots$ и $E=54.54825\ldots$ остаются верными для этого
квартичного завершения, но не являются уникальным следствием
канонического родителя.

\begin{proposition}[Уточнение без переписывания предыдущего гейта]
Предыдущий радиальный результат доказывает существование устойчивого
решения для одного допустимого положительного потенциала. Канонический
фазовый функционал проекта задаёт другое продолжение вне вакуумной орбиты и
должен проверяться отдельно.
\end{proposition}

\section{Канонический радиальный профиль}

Подстановка
\begin{equation}
 Q=\Delta u(r)\left(P-\frac{I_3}{3}\right),
 \qquad P=\widehat x\widehat x^T
\end{equation}
в \eqref{eq:v6-canonical-defect-potential} даёт барьер высоты
$0.0571782605\ldots$ между двумя вырожденными минимумами $u=0$ и $u=1$.
Краевая задача с прежними производными членами решена заново. Получены
\begin{align}
 r_{1/2}&=1.4343402812\ldots,\\
 E_D&=14.2011218910\ldots,\\
 E_F&=27.3834201076\ldots,\\
 E_V&=4.3940994196\ldots,\\
 E_{\rm can}&=45.9786414182\ldots.
 \label{eq:v6-canonical-radial-result}
\end{align}
Вириальный остаток равен $4.3\cdot10^{-8}$, а максимальный остаток
краевой задачи --- $2.0\cdot10^{-7}$.

Изотропный центр теперь не искусственно дорог: это вторая сосуществующая
фаза. Конечный радиус удерживается стенкой между фазами и положительной
кривизной составной связности.

\section{Пятикомпонентный гессиан}

Возмущение разлагалось в виде
\begin{equation}
 \delta Q(x)=\sum_{a,\ell,m,n}
 c_{a\ell mn}\,T_a\,S_{\ell m}(x)g_n(r),
 \label{eq:v6-five-component-galerkin-basis}
\end{equation}
где $T_a$ --- пять ортонормированных симметричных бесследовых матриц,
$S_{\ell m}$ --- вещественные регулярные гармонические полиномы при
$\ell=0,1,2$, а $g_n$ --- гладкие радиальные функции, обращающиеся в нуль
на границе шара.

Вторая вариация вычислена не заменой составной связности независимым полем,
а прямым дифференцированием
\begin{equation}
 A_i(Q)=\frac{[Q,\partial_iQ]}{\Delta^2},
 \qquad
 F_{ij}=\frac{2[\partial_iQ,\partial_jQ]}{\Delta^2}+[A_i,A_j].
 \label{eq:v6-composite-full-hessian-identities}
\end{equation}
Угловое интегрирование выполнено симметричной квадратурой, а радиальное ---
квадратурой Гаусса--Лежандра. Размеры последовательных пространств равны
$90,90,135,180$.

Нижние собственные значения обобщённого гессиана равны соответственно
\begin{equation}
 1.968249\ldots,
 \quad1.589985\ldots,
 \quad0.475785\ldots,
 \quad0.331040\ldots.
 \label{eq:v6-nonradial-lowest-sequence}
\end{equation}
Во всех четырёх пространствах отрицательных собственных значений не
найдено. Нижний уровень является почти трёхкратно вырожденным и лежит в
$\ell=2$-части выбранного скалярно-тензорного базиса.

Снижение значения при расширении базиса нельзя истолковывать как
возникновение отрицательной моды. У упорядоченного фона существуют
безмассовые директорные возбуждения, поэтому бесконечнообъёмный спектр не
обязан иметь положительную щель. Вычисление устанавливает только подход к
нулевому порогу сверху в проверенных пространствах.

\section{Независимые проверки гессиана}

Для двенадцати случайных направлений матричная квадратичная форма сравнена
с центральной второй разностью полного нелинейного функционала. Максимальный
относительный остаток равен $2.0\cdot10^{-7}$. Норма проекции первой
вариации не превышает $2.7\cdot10^{-5}$, а нарушение симметрии матрицы
гессиана машинно равно нулю.

\begin{proposition}[Конечномерный нерадиальный тест]
Канонический составной ёж не имеет отрицательной моды среди проверенных
пятикомпонентных возмущений с $\ell=0,1,2$ в пространствах Галёркина
размерности до $180$.
\end{proposition}

\section{Граница результата}

Получен сильный положительный вычислительный тест, но не теорема о полном
континуальном спектре. Не проверены произвольные радиальные функции,
гармоники $\ell\ge3$, существенно несферические конечные деформации и
нелинейное разделение дефекта на несколько объектов. Кроме того,
приближение нижней тройки к нулю требует отдельного различения
голдстоуновского порога, переносных направлений и возможного резонанса.

\begin{nogocriterion}[Конечный базис не доказывает полную устойчивость]
Отсутствие отрицательной моды в $180$-мерном пространстве запрещает
объявлять частицу окончательно устойчивой. Следующий гейт должен получить
радиальные матричные операторы фиксированного полного углового момента и
либо доказать их неотрицательность, либо найти канал распада.
\end{nogocriterion}

\begin{protocol}
\begin{enumerate}
 \item канонический объёмный потенциал: \textbf{восстановлен};
 \item равенство энергий изотропной и упорядоченной фаз:
 \textbf{учтено};
 \item прежнее квартичное завершение: \textbf{условно допустимо, но не
 уникально};
 \item каноническая радиальная энергия:
 \textbf{$45.9786414182\ldots$};
 \item пять компонент $Q$: \textbf{включены};
 \item угловые степени $\ell=0,1,2$: \textbf{проверены};
 \item максимальная размерность базиса: \textbf{$180$};
 \item отрицательные моды в проверенном базисе: \textbf{не найдены};
 \item положительная спектральная щель в бесконечном объёме:
 \textbf{не доказана};
 \item полная нелинейная устойчивость: \textbf{не доказана};
 \item следующий гейт: \textbf{континуальные канальные операторы и
 нулевой порог};
 \item устойчивая наблюдаемая частица: \textbf{ещё не получена}.
\end{enumerate}
\end{protocol}

Машинный сертификат сохранён в
\path{s2t_v6_bosonic_defect_nonradial_stability_gate_results.json}.